% ===================================================================
%  तालिका सं. 1 — विद्यालय। — माध्यमिक विद्यालय। — कक्षा।
%
%  Ceci n'est PAS une transcription : c'est une traduction, faite en
%  2026, en hindi d'aujourd'hui. D'ou trois differences avec texto/io
%  et texto/fr, et trois seulement :
%
%   * pas de \nl ni de \cc — la traduction ne suit aucune ligne
%     imprimee, et les lignes de ce fichier ne sont que du confort ;
%   * pas de \begin{VUpage} — elle n'a ni feuillet ni folio, et la
%     page de lecture ne lui en affiche pas ;
%   * le gras se pose sur le TERME seul, ponctuation dehors.
%
%  Tout le reste est identique, et doit l'etre : les cles « %%K » sont
%  celles de texto/io, macro pour macro pour l'apparat, et les renvois
%  \textsuperscript{(n)} visent les MEMES objets de la planche — ce
%  sont eux qui ouvrent les gros plans.
%
%  L'ido est le texte de base ; le francais sert de controle, et
%  l'anglais et l'espagnol de calibrage : les traductions doivent
%  s'accorder entre elles autant qu'avec la source.
%
%  LE GRAS NE SE POSE PAS SUR UN MOT-OUTIL. Le fac-simile ido le fait
%  par accident — « \VUgras{Ica} docisto » — et la page de lecture l'en
%  corrige ; la traduction, elle, n'a pas a repeter l'accident.
%
%  LES NOMS DE PERSONNE SONT NATURALISES, comme dans les six autres
%  colonnes : Ioannes est जॉन, Iakobus est जेम्स. Seuls le chien Medor
%  et l'ane Aliboron gardent leur nom.
%
%  LE HINDI POSTPOSE SES RAPPORTS et place le verbe a la fin : le
%  renvoi suit son substantif, comme partout, mais le substantif ne
%  se trouve pas toujours ou le francais le met. On garde l'ordre des
%  RENVOIS de l'ido, et l'on refait la phrase autour.
% ===================================================================

%%K t01-apar-0 apar
\VUsaut{2.0mm}
\VUcentre{12.6pt}{120}{व्याख्या-पुस्तिका}
\VUsaut{3.2mm}
\VUcentre{8.4pt}{160}{संलग्न}
\VUsaut{2.6mm}
\VUtitre{15.6pt}{0}{देल्मास सहायक तालिकाओं से}
\VUsaut{3.4mm}
\VUornamento{9pt}
\VUsaut{5.0mm}
\VUcentre{13.2pt}{90}{पहली शृंखला}
\VUsaut{3.0mm}
\VUfilet{16mm}
\VUsaut{5.6mm}

%%K t01-apar-1 apar
\VUtitre{15.0pt}{60}{तालिका सं. 1}
\VUsaut{1.4mm}
\VUtitre{11.4pt}{0}{विद्यालय। --- माध्यमिक विद्यालय। --- कक्षा।}
\VUsaut{2.2mm}
\VUfilet{20mm}
\VUsaut{6.4mm}

%%K t01-tit-1 sub
\VUpk{10.2pt}{40}{सामान्य विवरण।}
\VUsaut{3.0mm}

%%K t01-01-1 p
1. --- इस तालिका में एक \VUgras{माध्यमिक विद्यालय} की \VUgras{कक्षा} दिखाई गई
है। इस कक्षा में आठ \VUgras{मेज़ें} \textsuperscript{(18)} और आठ \VUgras{बेंच}
\textsuperscript{(32)} हैं। हर बेंच पर तीन विद्यार्थी बैठते हैं। \VUgras{अध्यापक}
\textsuperscript{(7)} \VUgras{कुर्सी} \textsuperscript{(8)} पर अपने \VUgras{मंच}
\textsuperscript{(6)} के पीछे बैठे हैं।

%%K t01-02-1 p
2. --- मंच के बाईं ओर काँच की एक \VUgras{अलमारी} \textsuperscript{(15)} खड़ी है।
दाईं ओर \VUgras{श्यामपट} \textsuperscript{(1)} है, और उस पर \VUgras{झाड़न}
\textsuperscript{(2)} तथा \VUgras{खड़िया} \textsuperscript{(3)}।

%%K t01-02-2 p
मंच के ऊपर \VUgras{भित्ति-चित्र} \textsuperscript{(9, 11, 12)} और एक
\VUgras{मानचित्र} \textsuperscript{(10)} टँगे हैं।

%%K t01-03-1 p
3. --- सब विद्यार्थी अपनी \VUgras{जगह} \textsuperscript{(17)} पर नहीं हैं। जॉन
\textsuperscript{(4)} श्यामपट के पास खड़ा है; वह झाड़न पकड़े हुए है और
\VUgras{ज्यामितीय आकृतियाँ} मिटा रहा है।

%%K t01-03-2 p
पीटर मंच के पास है; वह \VUgras{लिखित अभ्यास} \textsuperscript{(5)} इकट्ठे करके
अध्यापक के पास ले जा रहा है।

%%K t01-04-1 p
4. --- यह अध्यापक \VUgras{आधुनिक भाषाओं} के अध्यापक हैं। वे विद्यार्थियों को
विदेशी भाषाएँ \VUgras{(जर्मन, अंग्रेज़ी, स्पेनी, इतालवी, रूसी} या
\VUgras{फ़्रेंच)} बोलना, पढ़ना और लिखना सिखाते हैं।

%%K t01-05-1 p
5. --- कक्षा बहुत बड़ी और बहुत उजली है। सिरे पर हवा और उजाला आने के लिए दो
\VUgras{रोशनदानों} \textsuperscript{(78)} वाली एक चौड़ी \VUgras{खिड़की} और एक
\VUgras{काँच का दरवाज़ा} है।

%%K t01-05-2 p
यह कक्षा ठंडी नहीं है। बीच में उसे गरम रखने के लिए एक बहुत अच्छी
\VUgras{अँगीठी} \textsuperscript{(46)} रखी है, अपनी \VUgras{नली}
\textsuperscript{(45)} समेत।

%%K t01-05-3 p
दीवार पर \VUgras{खूँटियाँ} \textsuperscript{(58)} जड़ी हैं; उन पर विद्यार्थियों की
टोपियाँ और कोट टँगे रहते हैं।

%%K t01-tit-2 sub
\VUsaut{5.2mm}
\VUpk{10.2pt}{40}{खेल का प्रांगण।}
\VUsaut{3.6mm}

%%K t01-06-1 p
6. --- \VUgras{घड़ी} \textsuperscript{(63)} नौ बजकर पाँच मिनट बता रही है।

%%K t01-06-2 p
\VUgras{अवकाश} \textsuperscript{(76)} अभी समाप्त हुआ है और पाठ फिर आरंभ हो रहा है।
अवकाश \VUgras{प्रांगण} \textsuperscript{(75)} में होता है। हमें कुछ विद्यार्थी खेलते
दिखाई देते हैं। एक \VUgras{सहायक अध्यापक} \textsuperscript{(74)} उन पर नज़र रखे हुए
है।

%%K t01-07-1 p
7. --- प्रांगण में हमें और भी कई लोग दिखाई देते हैं: \VUgras{निरीक्षक}
\textsuperscript{(73)}, \VUgras{प्रधानाचार्य} \textsuperscript{(71)} और \VUgras{अनुशासक}
\textsuperscript{(72)}।

%%K t01-07-2 p
निरीक्षक विद्यालयों का निरीक्षण करते हैं, प्रधानाचार्य विद्यालय चलाते हैं,
अनुशासक विद्यार्थियों की पढ़ाई पर नज़र रखते हैं।

%%K t01-07-3 p
चौथे व्यक्ति \VUgras{चपरासी} \textsuperscript{(77)} हैं। वे अनुपस्थित विद्यार्थियों
के नाम लिखने के लिए बगल में एक रजिस्टर दबाए हुए हैं।

%%K t01-08-1 p
8. --- प्रांगण के सिरे पर \VUgras{व्यायाम के उपकरण} \textsuperscript{(70)} और
\VUgras{हस्तकला} \textsuperscript{(69)} की कार्यशाला है।

%%K t01-08-2 p
तालिका के दाईं ओर \VUgras{संगीत} और \VUgras{चित्रकला} के \VUgras{कक्ष} मुश्किल
से दिखाई देते हैं।

%%K t01-noto-1 noto
\VUnotes{143.2mm}{%
(*) \textit{नामों} के लिए भी यही सुझाया जाता है कि उन्हें उनके लातीनी रूप में
लिखा जाए। असंख्य रूपों से बचने का यही एक उपाय है। फिर, \textit{Giovanni,}
\textit{Jean, Johann, Jan, Hans, John, Ivan} में से चुनें भी तो कैसे? क्या
नामों का एक अलग कोश बनाया जाए? लातीनी कर्ता-रूप ले लें और कहें:
\VUgras{Ioannes, Ioanna; Iulius, Iulia; Iakobus; Andreas; Lukas.}
(संपूर्ण व्याकरण)।%
}

%%K t01-tit-3 sub
\VUpk{10.2pt}{40}{विस्तृत विवरण।}
\VUsaut{2.4mm}

%%K t01-tit-4 sub
\VUpk{10.2pt}{40}{अच्छे विद्यार्थी।}
\VUsaut{3.0mm}

%%K t01-09-1 p
9. --- मंच के दाईं ओर पहली बेंच पर बैठे विद्यार्थी \VUgras{हेनरी} (*),
\VUgras{स्टीफ़न} और \VUgras{चार्ल्स} हैं।

%%K t01-09-2 p
हेनरी बहुत ध्यान देने वाला और मेहनती है; वह अध्यापक की बात सुन रहा है। उसकी
मेज़ पर एक \VUgras{कापी} \textsuperscript{(28)}, एक \VUgras{पुस्तक}
\textsuperscript{(29)}, एक \VUgras{शब्दकोश} \textsuperscript{(30)} और एक
\VUgras{कलमदान} \textsuperscript{(31)} है। उसकी पुस्तक में बहुत से \VUgras{पृष्ठ}
हैं; हर अध्याय में कई \VUgras{अनुच्छेद} हैं और हर \VUgras{पंक्ति} में बहुत से
\VUgras{शब्द}।

%%K t01-09-4 p
उसका पड़ोसी स्टीफ़न \textsuperscript{(24)} उत्साही है। वह अगली बार का दिया हुआ पाठ
अपनी कापी में लिख रहा है। उसकी \VUgras{लिखावट} सुंदर है। उसकी पुस्तक बंद है।
केवल \VUgras{जिल्द} \textsuperscript{(27)} दिखाई देती है। उसके पास उसका
\VUgras{परकार} \textsuperscript{(26)} पड़ा है।

%%K t01-09-5 p
चार्ल्स \textsuperscript{(23)} हाथ में पुस्तक लिए है; वह उसे खोलकर अपना पाठ फिर
दोहरा रहा है। हर विद्यार्थी की तरह उसके सामने भी एक \VUgras{दवात}
\textsuperscript{(25)} है। वह आज्ञाकारी है।

%%K t01-10-1 p
10. --- अगली मेज़ पर \VUgras{लुई, एंटनी} और \VUgras{पॉल} हैं। लुई अपना
\VUgras{पाठ} \textsuperscript{(22)} सुना रहा है और अध्यापक के \VUgras{प्रश्नों} का
उत्तर दे रहा है। एंटनी \textsuperscript{(21)} बहुत असावधान है; अध्यापक कभी-कभी उसे
दंड भी देते हैं। पॉल \textsuperscript{(20)} अच्छा मित्र है, मिलनसार और उपकारी। वह
बहुत ध्यान देता है।

%%K t01-tit-5 sub
\VUsaut{4.2mm}
\VUpk{10.2pt}{40}{बुरे विद्यार्थी।}
\VUsaut{3.2mm}

%%K t01-11-1 p
11. --- हेनरी के पीछे \VUgras{जॉर्ज} \textsuperscript{(38)} बैठा है। वह बुरा लड़का
नहीं है, पर \VUgras{बातूनी}, असावधान और चंचल है। उसकी \VUgras{चपलता} और
\VUgras{अवज्ञा} पाठ के समय \VUgras{विघ्न} डालती हैं, और उसे प्रायः कड़ी
\VUgras{चेतावनी} मिलनी चाहिए। उसकी \VUgras{कच्ची कापी} \textsuperscript{(35)} गंदी
है, वह उस पर अपनी \VUgras{कलम} \textsuperscript{(34)} से \VUgras{स्याही के धब्बे}
डालता रहता है, इसीलिए उसे हमेशा अपना \VUgras{रबड़} \textsuperscript{(36)} चाहिए;
उसका \VUgras{बस्ता} \textsuperscript{(33)} फटा हुआ है, क्योंकि वह बहुत लापरवाह है।
आधुनिक भाषाओं की कक्षा में वह अपना \VUgras{पेंसिल-धारक} \textsuperscript{(37)} क्यों
लाता है?

%%K t01-11-2 p
उसका पड़ोसी एडमंड \textsuperscript{(39)} उसे पसंद नहीं करता। वह मानें तो थोड़ा
आलसी है, पर शांत और स्थिर है। वह बकबक नहीं करता; बस, सुनने के बदले वह अपनी
\VUgras{पाठ्य-पुस्तक} की \VUgras{कहानियाँ} पढ़ना अधिक पसंद करता है।

%%K t01-11-3 p
इस बेंच का तीसरा विद्यार्थी \VUgras{लुडोविक} \textsuperscript{(40)} है। वह परिश्रमी
है। वह \VUgras{कागज़ के} एक सफ़ेद \VUgras{पन्ने} पर लिख रहा है। अपने सामने उसने
अपना \VUgras{सोख़्ता} \textsuperscript{(41)} रखा है।

%%K t01-12-1 p
12. --- अध्यापक के बाईं ओर दूसरी बेंच के विद्यार्थी बहुत छोटे हैं। पहला, नन्हा
\VUgras{एमिल}, अभी अच्छी तरह लिख नहीं पाता; वह अपनी \VUgras{स्लेट}
\textsuperscript{(42)}, अपनी \VUgras{स्लेट-पेंसिल} और अपना \VUgras{स्पंज}
\textsuperscript{(43)} लाता है।

%%K t01-12-2 p
उसके पास \VUgras{ऑगस्ट} और \VUgras{बर्नार्ड} \textsuperscript{(44)} हैं। पिछला अच्छा
काम करता है, पर उसका \VUgras{वेश} बहुत मैला-कुचैला रहता है।

%%K t01-13-1 p
13. --- उनके पीछे तीन \VUgras{आलसी} बैठे हैं। \VUgras{अल्बर्ट} अपने पाठ याद नहीं
करता। \VUgras{जेम्स} \textsuperscript{(47)} सुनने के बदले सो रहा है। उसका
\VUgras{आलस्य} उसे \VUgras{डाँट} और \VUgras{दंड} तक दिलाता है। अंत में
\VUgras{क्रिश्चियन} \textsuperscript{(48)} कहीं अधिक चंचल है। उसका \VUgras{आचरण}
बुरा है और वह सुधरने का कोई यत्न नहीं करता।

%%K t01-14-1 p
14. --- उनके पड़ोसी, इसके विपरीत, सुशील हैं। \VUgras{अर्नेस्ट}
\textsuperscript{(51)} को क्रम और नियम प्रिय हैं; वह सदा अपना \VUgras{रूलर}
\textsuperscript{(50)} और अपनी \VUgras{पेंसिल} \textsuperscript{(49)} काम में लाता है।
वह \VUgras{बाह्य छात्र} है।

%%K t01-14-2 p
\VUgras{फ़्रांसिस} \textsuperscript{(52)} \VUgras{छात्रावासी} है; वह वरदी पहनता है,
और विद्यालय में ही खाता और सोता है। वह अच्छा \VUgras{मित्र} है।

%%K t01-14-3 p
मॉरिस अपने \VUgras{चाकू} \textsuperscript{(56)} से अपनी \VUgras{पेंसिल} छील रहा है;
वह बहुत सावधान है।

%%K t01-14-4 p
वह अपनी सब पुस्तकें लाता है, अपना \VUgras{व्याकरण} \textsuperscript{(55)}, अपनी
\VUgras{नोटबुक} \textsuperscript{(54)}, और एक \VUgras{लिखने का गद्दा}
\textsuperscript{(53)} भी। उसका \VUgras{थैला} \textsuperscript{(57)} उसके पीछे ज़मीन पर
है।

%%K t01-14-5 p
कक्षा के पीछे की दो मेज़ों पर \VUgras{अच्छे विद्यार्थी} \textsuperscript{(19)} बैठे
हैं।

%%K t01-tit-6 sub
\VUsaut{4.6mm}
\VUpk{10.2pt}{40}{चित्रकला और संगीत।}
\VUsaut{3.4mm}

%%K t01-15-1 p
15. --- \VUgras{चित्रकला-कक्ष} में दीवार पर सुंदर \VUgras{नमूने} टँगे हैं और
छत से एक \VUgras{छाया-आवरण} समेत \VUgras{लैंप} \textsuperscript{(62)} लटका है।
\VUgras{अध्यापक} \textsuperscript{(60)} बहुत धैर्यवान हैं; वे विद्यार्थियों के
\VUgras{चित्र} \textsuperscript{(61)} सुधारते हैं और उन्हें सामने रखी
\VUgras{प्रतिमा} \textsuperscript{(59)} देखकर चित्र बनाना सिखाते हैं।

%%K t01-16-1 p
16. --- \VUgras{संगीत-कक्ष} बहुत रोचक लगता है। वहाँ \VUgras{संगीत} पढ़ना,
\VUgras{स्वरलिपि} पर लिखे \VUgras{स्वर} \textsuperscript{(67)} गाना (डो, रे, मी,
फ़ा, सोल, ला, सी\,=\,C. D. E. F. G. A. B.), \VUgras{कुंजियाँ} पहचानना और
मनोहर \VUgras{धुनें} गाना सिखाया जाता है। स्वरलिपियाँ \VUgras{रिहल}
\textsuperscript{(65)} पर रखी जाती हैं। एक विद्यार्थी \VUgras{पियानो बजा रहा}
\textsuperscript{(64)} है और \VUgras{संगीत-अध्यापक} \textsuperscript{(66)}
\VUgras{ताल दे रहे} \textsuperscript{(68)} हैं।

%%K t01-17-1 p
17. --- हमें \VUgras{हस्तकला} \textsuperscript{(69)} की \VUgras{कार्यशाला} का एक
कोना मुश्किल से दिखाई देता है, जहाँ बच्चे लकड़ी छीलना और लोहा रेतना सीख सकते
हैं। यह हर विद्यालय में नहीं होती।

%%K t01-tit-7 sub
\VUsaut{5.0mm}
\VUpk{10.2pt}{40}{विज्ञान-कक्ष।}
\VUsaut{3.6mm}

%%K t01-18-1 p
18. --- गणित के अध्यापक विद्यार्थियों को \VUgras{ज्यामिति, अंकगणित, गणना} और
\VUgras{मीटरी प्रणाली} पढ़ाते हैं। तालिका पर हमें ज्यामिति की मुख्य आकृतियाँ
दिखाई देती हैं: एक \VUgras{वृत्त} \textsuperscript{(a)}, एक \VUgras{वर्ग}
\textsuperscript{(b)}, एक \VUgras{त्रिभुज} \textsuperscript{(c)}, एक \VUgras{रेखा}
\textsuperscript{(d)}।

%%K t01-18-2 p
हमें एक \VUgras{हिसाब} भी दिखाई देता है; वह न \VUgras{जोड़} है, न
\VUgras{घटाव}, न \VUgras{गुणा}: वह \VUgras{भाग} \textsuperscript{(e)} है।

%%K t01-19-1 p
19. --- मीटरी प्रणाली की तालिका पर हमें दिखाई देते हैं: एक \VUgras{मीटर}
\textsuperscript{(a)}, एक \VUgras{तराज़ू} \textsuperscript{(b)} और उसके \VUgras{बाट},
एक \VUgras{लीटर} \textsuperscript{(e)}, एक \VUgras{डेकालीटर} \textsuperscript{(f)}, एक
\VUgras{डेसीलीटर} और एक \VUgras{घन मीटर} \textsuperscript{(d)}।

%%K t01-19-2 p
विद्यार्थी \VUgras{प्राकृतिक विज्ञान} \textsuperscript{(11)} --- प्राणिविज्ञान
\textsuperscript{(a)}, वनस्पतिविज्ञान \textsuperscript{(b)}, भूविज्ञान
\textsuperscript{(c)} --- और \VUgras{इतिहास} \textsuperscript{(12)} भी पढ़ते हैं।

%%K t01-19-3 p
अलमारी में \VUgras{भौतिकी} \textsuperscript{(15)} और \VUgras{रसायन}
\textsuperscript{(16)} के उपकरण रखे हैं।

%%K t01-19-4 p
अलमारी के ऊपर हैं: एक \VUgras{गोला} \textsuperscript{(14)}, एक \VUgras{शंकु} और
एक \VUgras{भूगोलक} \textsuperscript{(13)}।

%%K t01-19-5 p
तालिकाओं के ऊपर एक \VUgras{मानचित्र} टँगा है जिस पर संसार के पाँचों भाग दिखाए
गए हैं: \VUgras{यूरोप} \textsuperscript{(g)}, \VUgras{एशिया} \textsuperscript{(e)},
\VUgras{अफ़्रीका} \textsuperscript{(f)}, \VUgras{अमेरिका} \textsuperscript{(ab)} और
\VUgras{ओशिआनिया} \textsuperscript{(h)}, तथा दो बड़े महासागर, \VUgras{प्रशांत}
\textsuperscript{(c)} और \VUgras{अटलांटिक} \textsuperscript{(d)}।
\VUsaut{6.0mm}
\VUfilet{22mm}
